\section*{Введение}
\addcontentsline{toc}{section}{Введение}

В главе~3 группа Галуа $G_f$ сепарабельного многочлена $f$ была определена как группа автоморфизмов его поля разложения
(\rf{3.27a}) и описана как группа перестановок корней, сохраняющих все полиномиальные соотношения между ними
(\rf{3.27b}), а основная теорема превратила вопросы о полях в вопросы о подгруппах. Теперь мы учимся эту группу
\emph{вычислять}. Инструментов несколько, и все они так или иначе сводят задачу к конечному числу проверок:
транзитивность действия (неприводимость $f$), дискриминант (лежит ли $G_f$ в знакопеременной группе), резольвенты (как
$G_f$ переставляет некоторые выражения от корней), редукция по модулю простых чисел (какие цикловые типы встречаются в
$G_f$). Для многочленов малых степеней этих сведений хватает, чтобы определить группу полностью.

\begin{itemize}[leftmargin=*]
\item \emph{Когда $G_f\subset A_n$?} (\rf{4.1}--\rf{4.4}). Дискриминант $D(f)=\prod_{i<j}(\alpha_i-\alpha_j)^2$ и квадратный
корень из него $\Delta(f)$, меняющий знак под действием нечётных перестановок (\rf{4.1}); при $\operatorname{char}F\ne2$
подполе, отвечающее $G_f\cap A_n$, есть $F[\sqrt{D(f)}]$, и $G_f\subset A_n$ тогда и только тогда, когда $D(f)$~--- квадрат
в $F$ (\rf{4.2}). Формулы для дискриминантов многочленов степеней $2$ и $3$, поведение при сдвиге и растяжении, принцип
специализации, позволяющий выводить такие формулы <<в общем положении>> (\rf{4.2a}--\rf{4.2c}); дискриминант через
производную и через критические значения (\rf{4.2d}), пять выводов формулы $D(X^3+pX+q)=-4p^3-27q^2$ (\rf{4.2e}) и
дискриминанты квартик (\rf{4.2f}). Знак дискриминанта над $\R$ равен $(-1)^s$, где $2s$~--- число невещественных корней, и
комплексное сопряжение как элемент $G_f$ (\rf{4.3}); вещественные кубики (\rf{4.3a}). В характеристике $2$ дискриминант
всегда квадрат, и его заменяет дискриминант Берлекэмпа: вместо извлечения квадратного корня решается уравнение
Артина--Шрайера (\rf{4.4}).
\item \emph{Транзитивность} (\rf{4.5}). Сепарабельный многочлен неприводим тогда и только тогда, когда $G_f$ действует на
его корнях транзитивно; тогда $G_f$~--- транзитивная подгруппа $S_n$ порядка, делящегося на $n=\deg f$.
\item \emph{Многочлены степени не выше трёх} (\rf{4.6}--\rf{4.7i}). Квадратные многочлены и особенности характеристики $2$
(\rf{4.6}); для неприводимой кубики $G_f=A_3$ или $S_3$ в зависимости от того, квадрат ли дискриминант, и сводка для всех
кубик, многочлены Шенкса (\rf{4.7}, \rf{4.7a}); формула Кардано: почему она даёт все корни и как согласовать
кубические корни (\rf{4.7b}, \rf{4.7c}; ср.\ \rf{3.28a}), резольвенты Лагранжа и башня Кардано (\rf{4.7d}, \rf{4.7e}),
примеры Кардано и Бомбелли (\rf{4.7f}), тригонометрическое решение в неприводимом случае (\rf{4.7g}, \rf{4.7h}),
преобразование Чирнгауза (\rf{4.7i}).
\item \emph{Многочлены четвёртой степени} (\rf{4.7j}--\rf{4.14d}). Три разбиения корней на пары (\rf{4.7j}), нормальная
подгруппа $V\subset S_4$ и неподвижное поле подгруппы $G_f\cap V$, кубическая резольвента и равенство её дискриминанта дискриминанту
квартики, явные формулы для дискриминанта квартики (\rf{4.8}--\rf{4.9a}); транзитивные подгруппы $S_4$ и как искать
транзитивные подгруппы $S_n$ (\rf{4.9b}, \rf{4.9c}); таблица возможных групп $S_4$, $A_4$, $V$, $D_4$, $C_4$ (\rf{4.9d}) и
способ различить $D_4$ и $C_4$~--- критерий Каппе--Уоррена (\rf{4.9e}--\rf{4.9g}). Решение уравнения четвёртой степени:
методы Феррари и Декарта, критерий разложимости на квадратичные множители, формула Эйлера, пучок коник, число вещественных
корней (\rf{4.9h}--\rf{4.9n}); примеры (\rf{4.10}--\rf{4.13c}) и задача о скрещённых лестницах (\rf{4.14}--\rf{4.14d}).
\item \emph{Многочлены с группой $S_p$ над $\Q$} (\rf{4.15}--\rf{4.18i}). Группа $S_p$ порождается любой транспозицией и
любым $p$-циклом (\rf{4.15}); неприводимый многочлен простой степени $p$ с ровно двумя невещественными корнями имеет
группу $S_p$ (\rf{4.16}, \rf{4.16a}); построение таких многочленов для каждого простого $p$ (\rf{4.16b}--\rf{4.17b}) и пример
того, что невещественные корни для этого не обязательны (\rf{4.18}). Транзитивные подгруппы $S_5$ и их цикловые типы
(\rf{4.18a}, \rf{4.18b}); метод резольвент и резольвента Кэли степени $6$, которая вместе с дискриминантом определяет
группу неприводимой квинтики (с точностью до выбора между $D_5$ и $C_5$; \rf{4.18c}--\rf{4.18g}); уравнение пятой
степени и модулярные функции (\rf{4.18h}); явное решение в радикалах квинтики с разрешимой группой (\rf{4.18i}).
\item \emph{Конечные поля} (\rf{4.19}--\rf{4.25a}). Расширения конечных полей просты (\rf{4.19}); поле $\F_q$ из $q=p^n$
элементов существует, единственно и является циклическим расширением $\F_p$, порождённым автоморфизмом Фробениуса
(\rf{4.20}); его подполя (\rf{4.21}); разложение $X^{p^n}-X$ на неприводимые множители и число неприводимых многочленов
данной степени (\rf{4.22}--\rf{4.22d}); алгебраическое замыкание поля $\F_p$ как объединение полей $\F_{p^n}$
(\rf{4.23}--\rf{4.25}); группа Галуа многочлена над конечным полем циклична и порождена Фробениусом, цикловой тип которого
задают степени неприводимых множителей (\rf{4.25a}).
\item \emph{Вычисление групп Галуа над $\Q$} (\rf{4.26}--\rf{4.34}). Орбиты $G_f$ и неприводимые множители $f$ (\rf{4.26},
\rf{4.27}); теорема Дедекинда с полным доказательством: если $f\in\Z[X]$ унитарен и $f\bmod p$ не имеет кратных корней,
то разложение $f$ по модулю $p$ на неприводимые множители степеней $m_1,\dots,m_r$ указывает элемент $G_f$ циклового типа
$(m_1,\dots,m_r)$ (\rf{4.27b}--\rf{4.28a}); элементы и классы Фробениуса (\rf{4.28b}--\rf{4.29}), теорема Чеботарёва о
плотности (\rf{4.29a}) и как угадать группу по частотам цикловых типов (\rf{4.29c}); примеры и построение многочленов с
группой $S_n$ (\rf{4.30}--\rf{4.32a}); общий алгоритм (\rf{4.32b}); <<почти все>> многочлены имеют группу $S_n$
(\rf{4.33}, \rf{4.34}).
\item \emph{Приложение: дискриминант через результант} (\rf{4.35}--\rf{4.39}). Результант и его свойства, позволяющие
вычислять его алгоритмом Евклида (\rf{4.35}); $D(f)=\pm\operatorname{Res}(f,f')$ и тонкость в характеристике, делящей
$\deg f$ (\rf{4.36}); дискриминанты $X^3+bX+c$, $X^5+X+1$ и $X^n+aX+b$ (\rf{4.37}--\rf{4.39}). Приложение можно читать
сразу после первого раздела; на его формулы мы ссылаемся уже там.
\item \emph{Упражнения} 4-1--4-10: поле разложения $X^m-1$ и вообще любого многочлена над $\F_p$ (\rf{4.39k}), группы Галуа
конкретных многочленов, поле разложения $X^8-2$~--- его группа порядка $16$ и решётка подполей (\rf{4.39l}), расширения
степени $4$ без промежуточных полей (\rf{4.39m}, \rf{4.39n}), неприводимые кубики над $\F_7$, многочлены $X^5+aX+b$ с
группой $D_5$; к упражнениям 4-8--4-10 даны указания.
\end{itemize}

Всюду в этой главе $F$~--- поле; для сепарабельного многочлена $f\in F[X]$ через $F_f$ обозначается его поле разложения,
а через $G_f=\Gal(F_f/F)$~--- его группа Галуа, которую мы отождествляем с подгруппой группы перестановок корней $f$
(\rf{3.27a}).
